% GENERATED, edit the French unit or the translation table
% GENERATED from content/data/enrichment.json by mxbuild.enrichment; do not edit.
\mxunit{enrichment-bayes-notes}{Reversing a conditional probability}
\mxsubtitle{High school · probability \textperiodcentered{} Lesson and method \textperiodcentered{} 07}

\begin{prerequisites}
Probability trees, intersections and conditional probability.
\end{prerequisites}

\begin{keypoint}[Ideas and method]
$P(A|B)$ is the probability of A given that B occurred. If $P(B) > 0$, $P(A|B) = \frac{P(A \cap  B)}{P(B)}$. The information B restricts the sample space.

$P(A \cap  B) = P(A)P(B|A)$. If A and its complement form the first two branches of a tree, $P(B) = P(A)P(B|A) + P(\text{not }A)P(B|\text{not }A)$.

Bayes’ formula follows from these two identities. $P(A|B)$ generally differs from $P(B|A)$: the denominators refer to different populations.
\end{keypoint}

\begin{example}
Choose urn A with probability $\frac{1}{3}$ and B with probability $\frac{2}{3}$. A contains 3 red and 1 blue ball; B contains 1 red and 3 blue balls. $P(\text{red}) = (\frac{1}{3})(\frac{3}{4}) + (\frac{2}{3})(\frac{1}{4}) = \frac{5}{12}$. Given a red ball, $P(A|\text{red}) = \frac{\frac{1}{4}}{\frac{5}{12}} = \frac{3}{5}$.
\end{example}

\begin{figurebox}[Only branches leading to a red ball are shown.]
\begin{tikzpicture}[x=1.6cm,y=.8cm]\draw[mxmint,thick](0,0)--(1,1)node[midway,above]{$\frac13$};\draw[mxmint,thick](0,0)--(1,-1)node[midway,below]{$\frac23$};\node at(1.2,1){$A$};\node at(1.2,-1){$B$};\draw[mxink,thick](1.5,1)--(3,1.5)node[midway,above]{$\frac34$};\draw[mxink,thick](1.5,-1)--(3,-.5)node[midway,above]{$\frac14$};\node[right]at(3,1.5){$A\cap R:\ \frac14$};\node[right]at(3,-.5){$B\cap R:\ \frac16$};\end{tikzpicture}
\end{figurebox}

\begin{activity}[Explore]
Represent the example with 120 hypothetical trials matching the proportions exactly. Count red balls from each urn.
\end{activity}

